\def\handoutsay{false}\def\handoutfill{false}
\documentclass[12pt]{article}
\newcommand{\handoutrunning}{Lecture 10: Normal Vectors and Quadratic Forms}
\providecommand{\handoutsay}{true}
\usepackage[T1]{fontenc}
\usepackage{lmodern,microtype}
\usepackage[a4paper,margin=15mm,headheight=14pt,headsep=5mm]{geometry}
\usepackage{amsmath,amssymb,amsthm,mathtools}
\usepackage{enumitem,booktabs,needspace}
\usepackage{pifont}
\usepackage[dvipsnames]{xcolor}
\usepackage{fancyhdr}
\usepackage{ifthen}
\usepackage[hidelinks]{hyperref}
\hypersetup{pdfauthor={Chang-Ye Tu}}
\pagestyle{fancy}
\fancyhf{}
\fancyhead[L]{\small Linear Algebra}
\fancyhead[R]{\small\handoutrunning}
\fancyfoot[C]{\small\thepage}
\renewcommand{\headrulewidth}{0.3pt}
\setlength{\parindent}{0pt}
\setlength{\parskip}{4pt plus 1pt minus 1pt}
\setlength{\emergencystretch}{1.5em}
\usepackage{etoolbox}
\newcommand{\handoutdisplayskips}{\setlength{\abovedisplayskip}{7pt plus 2pt minus 2pt}%
  \setlength{\belowdisplayskip}{7pt plus 2pt minus 2pt}}
\apptocmd{\normalsize}{\handoutdisplayskips}{}{\PackageWarning{handout}{could not patch normalsize}}
\handoutdisplayskips
\setlist[enumerate]{label=(\alph*),leftmargin=*,itemsep=2pt,topsep=3pt}
\setlist[itemize]{leftmargin=*,itemsep=2pt,topsep=3pt}
\newcommand{\R}{\mathbb R}
\newcommand{\C}{\mathbb C}
\newcommand{\iu}{\mathrm{i}}
\newcommand{\contradiction}{\mathord{\text{\large\ding{56}}}}
\newcommand{\one}{\mathbf 1}
\DeclareMathOperator{\tr}{tr}
\DeclareMathOperator{\col}{col}
\DeclareMathOperator{\diag}{diag}
\DeclareMathOperator{\rank}{rank}
\DeclareMathOperator\prb{\mathsf{P}}
\DeclareMathOperator\expc{\mathsf{E}}
\DeclareMathOperator\var{var}
\DeclareMathOperator\cov{cov}
\DeclareMathOperator\cor{cor}
\DeclareMathOperator\argmin{argmin}
\DeclareMathOperator\argmax{argmax}
\DeclareMathOperator{\Span}{span}
\theoremstyle{definition}
\newtheorem{theorem}{Theorem}
\newtheorem{proposition}[theorem]{Proposition}
\newtheorem{lemma}[theorem]{Lemma}
\newtheorem{corollary}[theorem]{Corollary}
\newtheorem{definition}[theorem]{Definition}
\newtheorem{example}[theorem]{Example}
\newtheorem{remark}[theorem]{Remark}
\newtheorem{fact}{Fact}
\renewcommand{\thefact}{F\arabic{fact}}
\AddToHook{cmd/subsection/before}{\par\Needspace{9\baselineskip}}
\AddToHook{env/theorem/before}{\par\Needspace{4\baselineskip}}
\AddToHook{env/proposition/before}{\par\Needspace{4\baselineskip}}
\AddToHook{env/lemma/before}{\par\Needspace{4\baselineskip}}
\AddToHook{env/definition/before}{\par\Needspace{4\baselineskip}}
\AddToHook{env/example/before}{\par\Needspace{4\baselineskip}}
\definecolor{saybg}{gray}{0.86}
\NewDocumentEnvironment{say}{+b}
  {\ifthenelse{\equal{\handoutsay}{true}}{\par\Needspace{2\baselineskip}\noindent\colorbox{saybg}{\parbox{\dimexpr\linewidth-2\fboxsep\relax}%
   {\small\sffamily\ding{43}\enspace\textbf{Say.}\enspace#1}}\par}{}}
  {}
\providecommand{\handoutfill}{false}
\newsavebox{\boardbox}
\newlength{\boardht}
\definecolor{boardrule}{gray}{0.55}
\makeatletter
\newcommand{\boardstart}{\par\if@inlabel\leavevmode\par\fi}
\makeatother
\NewDocumentEnvironment{board}{o +b}
  {\boardstart\IfValueT{#1}{\noindent#1\par\nopagebreak}%
   \noindent\begin{lrbox}{\boardbox}\begin{minipage}{\linewidth}#2\end{minipage}\end{lrbox}%
   \setlength{\boardht}{\dimexpr\ht\boardbox+\dp\boardbox\relax}%
   \ifthenelse{\equal{\handoutfill}{true}}%
     {\setlength{\boardht}{1.4\boardht}%
      {\setlength{\fboxsep}{0pt}\setlength{\fboxrule}{0.4pt}\color{boardrule}%
       \framebox[\dimexpr\linewidth-2\fboxrule\relax]{\rule{0pt}{\boardht}}}}%
     {\ifthenelse{\equal{\handoutsay}{true}}%
        {\llap{\color{boardrule}\rule[-\dp\boardbox]{1.5pt}{\boardht}\hspace{4pt}}}{}%
      \usebox{\boardbox}}%
   \par}
  {}
\newcommand{\blank}[1]{\ifthenelse{\equal{\handoutfill}{true}}{\underline{\phantom{#1}}}{#1}}
\newcommand{\lec}[1]{Lecture~#1}
\newcommand{\handouttitle}[2]{%
  \begin{center}
  {\LARGE\sffamily\bfseries #2\par}\vspace{3pt}
  {\sffamily Linear Algebra \textperiodcentered\ \lec{#1}\par}
  \ifthenelse{\equal{\handoutsay}{true}}{\vspace{3pt}{\small\sffamily\colorbox{saybg}{\ding{43}\enspace\textbf{Say.}}\enspace
    teacher's copy: the shaded boxes are not in the students' copy.\par}}{}
  \ifthenelse{\equal{\handoutfill}{true}}{\vspace{3pt}{\small\sffamily Fill-in copy: complete the framed spaces in
    class.\par}}{}
  \end{center}}

\hypersetup{pdftitle={Linear Algebra, Lecture 10: Normal Vectors and Quadratic Forms}}
\newcommand{\boardtop}{\vspace*{-\abovedisplayskip}}
\begin{document}
\handouttitle{10}{Normal Vectors and Quadratic Forms}

\textbf{Last time} (\lec{9}): \(\expc(Az+b)=A\mu+b\), \(\var(Az+b)=A\Sigma A'\), \(\cov(Az,Bz)=A\Sigma B'\); the
\(4\times4\) \(\Sigma\) of Example~5. F1--F5: Chapters 1--6 of the textbook; F6--F11: probability, no proof.
\(I=I_n\), \(e_j=\col_j(I)\).

\subsection*{Facts}

\begin{enumerate}[label=\textbf{F\arabic*.},ref=F\arabic*,leftmargin=*,itemsep=0pt,parsep=0pt,topsep=2pt]
\item\label{f-rules} \((AB)C=A(BC)\), \(A(B+C)=AB+AC\), \((A+B)C=AC+BC\), \(c(AB)=(cA)B=A(cB)\), \(I_mA=A=AI_n\)
(\(A\in\R^{m\times n}\)); \((A+B)'=A'+B'\), \((cA)'=cA'\), \((AB)'=B'A'\), \((A')'=A\); \(x'y=y'x\); \(A\) is
\textbf{symmetric} if \(A'=A\).
\item\label{f-views} \((AB)_{ij}=\sum_ra_{ir}b_{rj}\); for \(A=[a_1\ \cdots\ a_p]\): \(Ae_j=a_j\),
\(\col_j(AB)=A\col_j(B)\), \(AB=a_1b_1'+\dots+a_pb_p'\) (\(b_j'\) = row \(j\) of \(B\));
\(\diag(\alpha_1,\dots,\alpha_n)\diag(\beta_1,\dots,\beta_n)=\diag(\alpha_1\beta_1,\dots,\alpha_n\beta_n)\); block
matrices multiply, and are transposed, block by block (when the blocks fit).
\item\label{f-inverse} \(BA=I\) \(\Rightarrow\) \(B=A^{-1}\); \((A')^{-1}=(A^{-1})'\), so the inverse of an invertible
symmetric matrix is symmetric.
\item\label{f-centering} \(\one=\one_n=(1,\dots,1)'\in\R^n\), \(\bar y=\frac1n\one'y\), \(\bar J_n=\frac1n\one\one'\),
\(C_n=I_n-\bar J_n\): \(C_n'=C_n=C_n^2\), \(C_n\one=0\), \(C_ny=y-\bar y\one\), \(y'C_ny=\sum_i(y_i-\bar y)^2\).
\item\label{f-orthonormal} \(\lVert x\rVert^2=x'x\). \(Q\in\R^{n\times p}\) has \textbf{orthonormal} columns \(\iff\)
\(Q'Q=I_p\). If \(p=n\) (\(Q\) \textbf{orthogonal}): \(Q^{-1}=Q'\), \(QQ'=I_n\).
\item\label{f-prob} (Probability.) A random variable is a number whose value depends on chance. For random variables
\(u\), \(v\) and numbers \(a\), \(b\), \(c\): \(\expc(au+bv+c)=a\expc(u)+b\expc(v)+c\) (also for finite sums);
\(\var(u)=\expc\bigl((u-\expc(u))^2\bigr)\); \(\cov(u,v)=\expc\bigl((u-\expc(u))(v-\expc(v))\bigr)\). \(\prb(z\in A)\)
is the probability that a random vector \(z\in\R^n\) takes a value in \(A\subseteq\R^n\):
\(\prb(z\in A\cup B)=\prb(z\in A)+\prb(z\in B)\) for disjoint \(A\), \(B\), and \(\prb(z\in\R^n)=1\); \(y\), \(z\) have the
\textbf{same distribution} if \(\prb(y\in A)=\prb(z\in A)\) for all \(A\), and then \(\expc g(y)=\expc g(z)\).
\(u_1,\dots,u_m\) are \textbf{independent} if
\(\prb(u_1\in A_1,\dots,u_m\in A_m)=\prb(u_1\in A_1)\cdots\prb(u_m\in A_m)\) for all \(A_1,\dots,A_m\). If \(u\) takes
only the values \(a_1,\dots,a_m\): \(\expc g(u)=g(a_1)\prb(u=a_1)+\dots+g(a_m)\prb(u=a_m)\). (Sets and functions are
those for which the probabilities and expectations are defined.)
\item\label{f-factor} Random vectors \(u_1\in\R^{n_1},\dots,u_m\in\R^{n_m}\) whose MGFs (Definition~\ref{h-def-mvn})
are finite near \(0\) are independent \(\iff\)
\(M_{(u_1',\dots,u_m')'}(t_1,\dots,t_m)=M_{u_1}(t_1)\cdots M_{u_m}(t_m)\) for all \(t_k\in\R^{n_k}\) near \(0\).
\item\label{f-functions} If \(u_1,\dots,u_m\) are independent, so are \(g_1(u_1),\dots,g_m(u_m)\), for any functions
\(g_1,\dots,g_m\).
\item\label{f-moments} If \(M_z\) is finite near \(0\), its first and second partial derivatives exist there, and at
\(t=0\): \(\partial M_z/\partial t_i=\expc(z_i)\), \(\partial^2M_z/\partial t_j\partial t_i=\expc(z_iz_j)\).
\item\label{f-product} If \(u_1,\dots,u_m\) are independent:
\(\expc\bigl(g_1(u_1)\cdots g_m(u_m)\bigr)=\expc g_1(u_1)\cdots\expc g_m(u_m)\) (finite expectations). So independent
random vectors \(u\), \(v\) have \(\cov(u,v)=0\).
\item\label{f-chisq} Let \(r,\nu,a,b\geqslant1\). \(\chi^2_r\) (\textbf{chi-square}) is the distribution of
\(w_1^2+\dots+w_r^2\) for independent \(w_1,\dots,w_r\sim N(0,1)\); \(t_\nu\) is that of \(Z/\sqrt{V/\nu}\) for
independent \(Z\sim N(0,1)\), \(V\sim\chi^2_\nu\); \(F_{a,b}\) is that of \((U/a)/(V/b)\) for independent
\(U\sim\chi^2_a\), \(V\sim\chi^2_b\). For \(u\) with one of these distributions, \(G(x)=\prb(u\leqslant x)\) is
continuous, so \(\prb(u=c)=0\) for every \(c\); \(G\) is strictly increasing on \(\R\) (\(t_\nu\)), or \(G(0)=0\) and
\(G\) is strictly increasing on \([0,\infty)\) (\(\chi^2_r\), \(F_{a,b}\)). For \(\tau\in(0,1)\), the \(\tau\)
\textbf{quantile} \(t_{\nu,\tau}\), \(\chi^2_{r,\tau}\), \(F_{a,b,\tau}\) is the unique \(c\) with \(G(c)=\tau\); so
\(\prb(u>c)=\alpha\) for the \((1-\alpha)\) quantile \(c\).
\end{enumerate}

\textbf{From \lec{8} and \lec{9}} (restated; \(A\in\R^{n\times n}\)).
\begin{itemize}[itemsep=0pt,parsep=0pt,topsep=2pt]
\item \lec{8}, Theorem~9 (\(A\) symmetric): (a) \(\R^n\) has an orthonormal basis \(q_1,\dots,q_n\) of eigenvectors of
\(A\). (b) If \(q_1,\dots,q_n\in\R^n\) are orthonormal, \(Aq_i=\lambda_iq_i\), \(Q=[q_1\ \cdots\ q_n]\) and
\(\Lambda=\diag(\lambda_1,\dots,\lambda_n)\): \(Q\) is orthogonal, \(A=Q\Lambda Q'\), the \(\lambda_i\) are the
eigenvalues of \(A\), and \(\rank A=\#\{i:\lambda_i\ne0\}\). (d) \(A=\lambda_1q_1q_1'+\dots+\lambda_nq_nq_n'\).
\item \lec{8}, Theorem~11: for symmetric \(A\), \(A^2=A\iff\) every eigenvalue of \(A\) is \(0\) or \(1\).
Example~5: \(\rank C_n=n-1\).
\item \lec{8}, Definition~13, Theorem~14(a),(c): a symmetric \(A\) is \textbf{positive definite} (PD) if
\(x'Ax>0\) for all \(x\ne0\), \textbf{positive semidefinite} (PSD) if \(x'Ax\geqslant0\) for all \(x\); with
\(\lambda_i\) as in Theorem~9(b): PD \(\iff\) every \(\lambda_i>0\); PSD \(\iff\) every \(\lambda_i\geqslant0\); PD
\(\Rightarrow\) invertible. Theorem~15(a): \(B\) PSD, \(K\in\R^{n\times k}\) \(\Rightarrow\) \(K'BK\) PSD.
\item \lec{9}, Definitions~1, 2: for a random vector \(z\), \(\expc(z)=(\expc(z_i))\),
\(\var(z)=\expc\bigl((z-\mu)(z-\mu)'\bigr)\) (\(\mu=\expc(z)\)) has the entries \(\sigma_{ij}=\cov(z_i,z_j)\);
\(\cov(u,v)=\expc\bigl((u-\mu_u)(v-\mu_v)'\bigr)\) has the entries \(\cov(u_i,v_j)\); so \(\var\bigl((u',v')'\bigr)\)
has the blocks \(\var(u)\), \(\cov(u,v)\); \(\cov(v,u)\), \(\var(v)\).
\item \lec{9}, Theorem~3(b),(c): \(\var(Az+b)=A\var(z)A'\); \(\var(z)\) is PSD. Theorem~4(d):
\(\expc(uu')=\var(u)+\mu_u\mu_u'\).
\item \lec{9}, Example~5: for the \(\Sigma\) of Example~\ref{h-exm-marginal} below, \(a=(-1,-2,1,0)'\) and
\(b=(2,2,-2,1)'\): \(\Sigma a=(0,0,1,2)'\), \(\Sigma b=(0,0,0,1)'\), \(a'\Sigma a=b'\Sigma b=1\),
\(\one'\Sigma\one=51\); \(\var(z_1+z_2)=\var(z_1-z_2)=5\), \(\cov(z_1+z_2,z_1-z_2)=3\) if \(\var(z)=\Sigma\).
\end{itemize}

\subsection*{1. Normal vectors}

\begin{definition}\label{h-def-mvn}
The \textbf{moment generating function} (MGF) of a random vector \(z\in\R^n\) is \(M_z(t)=\expc\exp(t'z)\)
(\(t\in\R^n\), where finite). Let \(\mu\in\R^n\) and \(\Sigma\in\R^{n\times n}\) be symmetric and PSD. \(z\) has the
\textbf{multivariate normal distribution} \(N_n(\mu,\Sigma)\), \(z\sim N_n(\mu,\Sigma)\), if
\[
M_z(t)=\exp\bigl(t'\mu+\tfrac12t'\Sigma t\bigr)\qquad\text{for all }t\in\R^n
\]
For \(n=1\): \(N(\mu,\sigma^2)\), with the MGF \(\exp(t\mu+\frac12\sigma^2t^2)\); \(N(0,1)\) is the \textbf{standard
normal} distribution. \(u\) and \(v\) are \textbf{jointly normal} if \((u',v')'\) is multivariate normal.
\end{definition}

\Needspace{6\baselineskip}
\begin{proposition}\label{h-prop-moments}
If \(z\sim N_n(\mu,\Sigma)\), then \(\expc(z)=\mu\) and \(\var(z)=\Sigma\).
\end{proposition}
\begin{proof}
The exponent \(q(t)=t'\mu+\frac12t'\Sigma t=\sum_kt_k\mu_k+\frac12\sum_k\sum_l\sigma_{kl}t_kt_l\) (\(\sigma_{kl}=\sigma_{lk}\))
and the chain and product rules give
\begin{align*}
\frac{\partial q}{\partial t_i}&=\mu_i+(\Sigma t)_i,\qquad\frac{\partial^2q}{\partial t_j\partial t_i}=\sigma_{ij},\qquad
\frac{\partial M_z}{\partial t_i}=\bigl(\mu_i+(\Sigma t)_i\bigr)M_z(t)\\
\frac{\partial^2M_z}{\partial t_j\partial t_i}&=\Bigl(\sigma_{ij}+\bigl(\mu_i+(\Sigma t)_i\bigr)\bigl(\mu_j+(\Sigma t)_j\bigr)\Bigr)M_z(t)
\end{align*}
At \(t=0\), \(M_z(0)=1\): \(\expc(z_i)=\mu_i\), \(\expc(z_iz_j)=\sigma_{ij}+\mu_i\mu_j\) (\ref{f-moments}). So
\(\expc(zz')=\Sigma+\mu\mu'\) and \(\var(z)=\Sigma\) (\lec{9}, Theorem~4(d)).
\end{proof}

\Needspace{10\baselineskip}
\begin{theorem}\label{h-thm-linear}
Let \(z\sim N_n(\mu,\Sigma)\), and \(A\in\R^{m\times n}\), \(b\in\R^m\) constant. Then
\(Az+b\sim N_m(A\mu+b,\ A\Sigma A')\).
\end{theorem}
\begin{proof}
\begin{board}[\(A\Sigma A'=(A')'\Sigma(A')\) is symmetric and PSD (\ref{f-rules}; \lec{8}, Theorem~15(a), \(K=A'\)); for
\(t\in\R^m\):]
\boardtop
\begin{align*}
t'(Az+b)&=t'b+(A't)'z\implies M_{Az+b}(t)=\exp(t'b)\,M_z(A't)&&\text{(\ref{f-rules}; \ref{f-prob})}\\
M_{Az+b}(t)&=\exp\bigl(t'b+(A't)'\mu+\tfrac12(A't)'\Sigma(A't)\bigr)&&\text{(Definition~\ref{h-def-mvn})}\\
&=\exp\bigl(t'(A\mu+b)+\tfrac12t'(A\Sigma A')t\bigr)&&\text{(\ref{f-rules})}
\end{align*}
\end{board}
\nopagebreak
\end{proof}

\begin{corollary}\label{h-cor-marginal}
Let \(z\sim N_n(\mu,\Sigma)\).
\begin{enumerate}
\item For \(1\leqslant i_1<\dots<i_k\leqslant n\): \((z_{i_1},\dots,z_{i_k})'\sim N_k\), with the mean
\((\mu_{i_1},\dots,\mu_{i_k})'\) and the submatrix of \(\Sigma\) in the rows and columns \(i_1,\dots,i_k\) as covariance
matrix. In particular \(z_i\sim N(\mu_i,\sigma_{ii})\); and for \(z=(z_1',z_2')'\), \(z_1\in\R^p\), \(z_2\in\R^q\),
\(\Sigma_{11}\in\R^{p\times p}\),
\[
\mu=\begin{pmatrix}\mu_1\\\mu_2\end{pmatrix},\qquad\Sigma=\begin{pmatrix}\Sigma_{11}&\Sigma_{12}\\\Sigma_{21}&\Sigma_{22}\end{pmatrix}
\implies z_1\sim N_p(\mu_1,\Sigma_{11}),\quad z_2\sim N_q(\mu_2,\Sigma_{22})
\]
\item \(a'z\sim N(a'\mu,\ a'\Sigma a)\) for \(a\in\R^n\).
\end{enumerate}
\end{corollary}
\begin{proof}
(a) \(T=[e_{i_1}\ \cdots\ e_{i_k}]\): \(T'z=(z_{i_1},\dots,z_{i_k})'\), and \(T'\Sigma T\) has the entries
\(e_{i_r}'\Sigma e_{i_s}=\sigma_{i_ri_s}\) (\ref{f-views}); Theorem~\ref{h-thm-linear} with \(A=T'\), \(b=0\).

(b) Theorem~\ref{h-thm-linear} with \(A=a'\), \(b=0\).
\end{proof}

\Needspace{18\baselineskip}
\begin{example}\label{h-exm-marginal}
Let \(z\sim N_4(\mu,\Sigma)\), with \(\mu=(1,0,2,2)'\) and \(\Sigma\) of \lec{9}, Example~5:
\[
\Sigma=\begin{pmatrix}4&0&4&0\\0&1&2&2\\4&2&9&6\\0&2&6&9\end{pmatrix}=KK',\qquad
K=\begin{pmatrix}2&0&0&0\\0&1&0&0\\2&2&1&0\\0&2&2&1\end{pmatrix}
\]
(\(x'\Sigma x=\lVert K'x\rVert^2\geqslant0\): \(\Sigma\) is PSD.) Find the distributions of \(z_2\), \((z_1,z_3)'\), \(\one'z\), \(e=z_3-z_1-2z_2\), \(f=z_4+2z_1+2z_2-2z_3\) and
\((z_1+z_2,z_1-z_2)'\).
\end{example}
\textit{Solution.}
\begin{board}[Corollary~\ref{h-cor-marginal}; \(e=a'z\), \(f=b'z\), \(a'\Sigma a=b'\Sigma b=1\), \(\one'\Sigma\one=51\)
(\lec{9}, Example~5):]
\boardtop
\begin{align*}
z_2&\sim N(0,1),\qquad\begin{pmatrix}z_1\\z_3\end{pmatrix}\sim
N_2\left(\begin{pmatrix}1\\2\end{pmatrix},\begin{pmatrix}4&4\\4&9\end{pmatrix}\right)&&\text{(Corollary~\ref{h-cor-marginal}(a))}\\
\one'z&\sim N(5,51)&&(\text{(b)};\ \one'\mu=1+0+2+2=5)\\
e&\sim N(1,1)&&(a'\mu=-1-0+2+0=1)\\
f&\sim N(0,1)&&(b'\mu=2+0-4+2=0)
\end{align*}
\end{board}
\begin{board}[\(A=\left(\begin{smallmatrix}1&1&0&0\\1&-1&0&0\end{smallmatrix}\right)\), \(Az=(z_1+z_2,z_1-z_2)'\);
\(A\Sigma A'=\var(Az)\) (\lec{9}, Theorem~3(b)), with the entries of \lec{9}, Example~5:]
\[
A\mu=\begin{pmatrix}1+0\\1-0\end{pmatrix}=\begin{pmatrix}1\\1\end{pmatrix},\qquad
Az\sim N_2\left(\begin{pmatrix}1\\1\end{pmatrix},\begin{pmatrix}5&3\\3&5\end{pmatrix}\right)\qquad
\text{(Theorem~\ref{h-thm-linear})}
\]
\end{board}

\Needspace{22\baselineskip}
\subsection*{2. Independence}

\Needspace{19\baselineskip}
\begin{theorem}\label{h-thm-indep}
\begin{enumerate}
\item Let \(z=(u',v')'\sim N_{p+q}(\mu,\Sigma)\), \(u\in\R^p\), \(v\in\R^q\), \(\mu=(\mu_u',\mu_v')'\),
\(\Sigma=\left(\begin{smallmatrix}\Sigma_{uu}&\Sigma_{uv}\\\Sigma_{vu}&\Sigma_{vv}\end{smallmatrix}\right)\)
(\(\Sigma_{uu}\in\R^{p\times p}\)). Then \(u\), \(v\) are independent \(\iff\) \(\Sigma_{uv}=0\) (\(\iff\)
\(\cov(u,v)=0\)).
\item Let \(D=\diag(\sigma_1^2,\dots,\sigma_n^2)\). Then \(z\sim N_n(\mu,D)\) \(\iff\) \(z_1,\dots,z_n\) are independent
and \(z_i\sim N(\mu_i,\sigma_i^2)\). In particular \(z\sim N_n(\mu_0\one,\sigma^2I_n)\) \(\iff\) \(z_1,\dots,z_n\) are
independent \(N(\mu_0,\sigma^2)\) variables (a \textbf{random sample} from \(N(\mu_0,\sigma^2)\)).
\item Let \(z\sim N_n(\mu,\Sigma)\), and \(A\in\R^{p\times n}\), \(B\in\R^{q\times n}\) constant. Then \(Az\), \(Bz\) are
independent \(\iff\) \(A\Sigma B'=0\).
\end{enumerate}
\end{theorem}
\begin{proof}
(a) \(\Rightarrow\): \(\Sigma=\var(z)\) (Proposition~\ref{h-prop-moments}) has the block
\(\Sigma_{uv}=\cov(u,v)\) (\lec{9}, Definitions~1, 2), which is \(0\) if \(u\), \(v\) are independent (\ref{f-product}).
\begin{board}[(a) \(\Leftarrow\): \(\Sigma_{uv}=0\), so \(\Sigma_{vu}=\Sigma_{uv}'=0\); for \(s\in\R^p\), \(t\in\R^q\), by
blocks (\ref{f-views}):]
\boardtop
\begin{align*}
(s',t')\mu&=s'\mu_u+t'\mu_v,\qquad(s',t')\,\Sigma\begin{pmatrix}s\\t\end{pmatrix}=s'\Sigma_{uu}s+t'\Sigma_{vv}t\\
M_z(s,t)&=\exp\bigl(s'\mu_u+\tfrac12s'\Sigma_{uu}s\bigr)\exp\bigl(t'\mu_v+\tfrac12t'\Sigma_{vv}t\bigr)=M_u(s)M_v(t)
&&\text{(Corollary~\ref{h-cor-marginal}(a))}
\end{align*}
For all \(s\), \(t\): \(u\), \(v\) are independent (\ref{f-factor}).
\end{board}
\begin{board}[(b) \(t'Dt=\sum_i\sigma_i^2t_i^2\) (\ref{f-views}), so \(D\) is PSD; \(\Rightarrow\) uses
\(z_i\sim N(\mu_i,\sigma_i^2)\) (Corollary~\ref{h-cor-marginal}(a)), \(\Leftarrow\) uses \ref{f-product}; for all
\(t\in\R^n\):]
\boardtop
\begin{align*}
\Rightarrow\quad M_z(t)&=\exp\bigl(t'\mu+\tfrac12t'Dt\bigr)=\textstyle\prod_i\exp\bigl(t_i\mu_i+\tfrac12\sigma_i^2t_i^2\bigr)=\prod_iM_{z_i}(t_i)\\
\Leftarrow\quad M_z(t)&=\expc\textstyle\prod_ie^{t_iz_i}=\prod_i\expc e^{t_iz_i}=\prod_i\exp\bigl(t_i\mu_i+\tfrac12\sigma_i^2t_i^2\bigr)=\exp\bigl(t'\mu+\tfrac12t'Dt\bigr)
\end{align*}
\(\Rightarrow\): \(z_1,\dots,z_n\) are independent (\ref{f-factor}). \(\Leftarrow\): \(z\sim N_n(\mu,D)\)
(Definition~\ref{h-def-mvn}).
\end{board}
\begin{board}[(c) \(G=\left(\begin{smallmatrix}A\\B\end{smallmatrix}\right)\in\R^{(p+q)\times n}\),
\(Gz=\bigl((Az)',(Bz)'\bigr)'\) (\ref{f-views}):]
\[
Gz\sim N_{p+q}(G\mu,G\Sigma G'),\qquad G\Sigma G'=\begin{pmatrix}A\Sigma A'&A\Sigma B'\\B\Sigma A'&B\Sigma B'\end{pmatrix}
\qquad\text{(Theorem~\ref{h-thm-linear}; \ref{f-views})}
\]
By (a) with \(u=Az\), \(v=Bz\): independent \(\iff\) \(A\Sigma B'=0\).
\end{board}
\nopagebreak
\end{proof}

\begin{remark}\label{h-rem-uncorr}
Joint normality is needed in Theorem~\ref{h-thm-indep}(a). Let \(u\) take the values \(-1\), \(0\), \(1\), each with
probability \(\frac13\), and \(v=u^2\). By \ref{f-prob}, \(\expc(u)=\expc(u^3)=0\), \(\expc(u^2)=\frac23\), and
\(\cov(u,v)=\expc(u^3)-\expc(u)\expc(u^2)=0\); but \(\prb(u=0,\ v=1)=0\ne\frac13\cdot\frac23=\prb(u=0)\prb(v=1)\).
\end{remark}

\Needspace{11\baselineskip}
\begin{example}\label{h-exm-indep}
For \(z\) of Example~\ref{h-exm-marginal}: are \(z_1\) and \((z_2,z_4)'\) independent? \(e\) and \((z_1,z_2)'\)? \(f\)
and \((z_1,z_2,z_3)'\)? \(z_1+z_2\) and \(z_1-z_2\)?
\end{example}
\textit{Solution.}
\begin{board}[Corollary~\ref{h-cor-marginal}(a) with \(i_1,i_2,i_3=1,2,4\):]
\[
\begin{pmatrix}z_1\\z_2\\z_4\end{pmatrix}\sim
N_3\left(\begin{pmatrix}1\\0\\2\end{pmatrix},\begin{pmatrix}4&0&0\\0&1&2\\0&2&9\end{pmatrix}\right)
\]
\(\sigma_{12}=\sigma_{14}=0\): \(z_1\) and \((z_2,z_4)'\) are independent (Theorem~\ref{h-thm-indep}(a), \(p=1\),
\(q=2\)).
\end{board}
\begin{board}[Theorem~\ref{h-thm-indep}(c) with \(B=a'\), \(B=b'\); \(\Sigma a=(0,0,1,2)'\), \(\Sigma b=(0,0,0,1)'\)
(\lec{9}, Example~5):]
\boardtop
\begin{align*}
[I_2\ \ 0]\,\Sigma a&=(0,0)'\implies e\text{ and }(z_1,z_2)'\text{ independent},\qquad z_3=z_1+2z_2+e\\
[I_3\ \ 0]\,\Sigma b&=(0,0,0)'\implies f\text{ and }(z_1,z_2,z_3)'\text{ independent},\qquad z_4=-2z_1-2z_2+2z_3+f
\end{align*}
\(\cov(z_1+z_2,z_1-z_2)=3\ne0\): \(z_1+z_2\) and \(z_1-z_2\) are not independent (\ref{f-product}).
\end{board}

\Needspace{22\baselineskip}
\subsection*{3. Conditional distributions}

\Needspace{19\baselineskip}
\begin{theorem}\label{h-thm-conditional}
Let \(z\sim N_n(\mu,\Sigma)\) be partitioned as in Corollary~\ref{h-cor-marginal}(a), with \(\Sigma_{11}\) PD, and
\(\Sigma_{22\cdot1}=\Sigma_{22}-\Sigma_{21}\Sigma_{11}^{-1}\Sigma_{12}\). Then \(\Sigma_{22\cdot1}\) is PSD, and
\[
w=z_2-\mu_2-\Sigma_{21}\Sigma_{11}^{-1}(z_1-\mu_1)\sim N_q(0,\Sigma_{22\cdot1}),\qquad w\text{ and }z_1\text{ independent}
\]
So \(z_2=\mu_2+\Sigma_{21}\Sigma_{11}^{-1}(z_1-\mu_1)+w\). For \(a\in\R^p\),
\(N_q\bigl(\mu_2+\Sigma_{21}\Sigma_{11}^{-1}(a-\mu_1),\ \Sigma_{22\cdot1}\bigr)\) is the \textbf{conditional
distribution of \(z_2\) given \(z_1=a\)}.
\end{theorem}
\begin{proof}
\begin{board}[\(\Sigma_{11}^{-1}\) exists (\lec{8}, Theorem~14(c)) and is symmetric (\ref{f-inverse});
\(B=\Sigma_{21}\Sigma_{11}^{-1}\), \(B'=\Sigma_{11}^{-1}\Sigma_{12}\) (\(\Sigma_{21}'=\Sigma_{12}\));
\(L=\left(\begin{smallmatrix}I_p&0\\-B&I_q\end{smallmatrix}\right)\); blocks (\ref{f-views}):]
\boardtop
\begin{align*}
L(z-\mu)&=\begin{pmatrix}z_1-\mu_1\\-B(z_1-\mu_1)+z_2-\mu_2\end{pmatrix}=\begin{pmatrix}z_1-\mu_1\\w\end{pmatrix}\\
L\Sigma&=\begin{pmatrix}\Sigma_{11}&\Sigma_{12}\\\Sigma_{21}-B\Sigma_{11}&\Sigma_{22}-B\Sigma_{12}\end{pmatrix}
=\begin{pmatrix}\Sigma_{11}&\Sigma_{12}\\0&\Sigma_{22\cdot1}\end{pmatrix}\\
L\Sigma L'&=\begin{pmatrix}\Sigma_{11}&\Sigma_{12}\\0&\Sigma_{22\cdot1}\end{pmatrix}\begin{pmatrix}I_p&-B'\\0&I_q\end{pmatrix}
=\begin{pmatrix}\Sigma_{11}&\Sigma_{12}-\Sigma_{11}B'\\0&\Sigma_{22\cdot1}\end{pmatrix}
=\begin{pmatrix}\Sigma_{11}&0\\0&\Sigma_{22\cdot1}\end{pmatrix}
\end{align*}
\end{board}
\begin{board}[Theorem~\ref{h-thm-linear} with \(A=L\), \(b=-L\mu\):]
\boardtop
\begin{align*}
L(z-\mu)&\sim N_n(0,L\Sigma L')\implies w\sim N_q(0,\Sigma_{22\cdot1})&&\text{(Corollary~\ref{h-cor-marginal}(a))}\\
\Sigma_{22\cdot1}&=\var(w)\text{ is PSD}&&\text{(Proposition~\ref{h-prop-moments}; \lec{9}, Theorem~3(c))}
\end{align*}
The off-diagonal blocks of \(L\Sigma L'\) are \(0\): \(z_1-\mu_1\) and \(w\) are independent
(Theorem~\ref{h-thm-indep}(a)), and so are \(z_1=(z_1-\mu_1)+\mu_1\) and \(w\) (\ref{f-functions}).
\end{board}
\nopagebreak
\end{proof}
The conditional mean is an affine function of \(a\), and \(\Sigma_{22\cdot1}\) does not depend on \(a\). With
\(B=\Sigma_{21}\Sigma_{11}^{-1}\), \(B\Sigma_{11}B'=\Sigma_{21}\Sigma_{11}^{-1}\Sigma_{11}\Sigma_{11}^{-1}\Sigma_{12}
=\Sigma_{22}-\Sigma_{22\cdot1}\) is PSD (\lec{8}, Theorem~15(a) with \(\Sigma_{11}\), \(B'\) for its \(B\), \(K\)); so
\(\var(c'w)=c'\Sigma_{22\cdot1}c\leqslant c'\Sigma_{22}c=\var(c'z_2)\) for \(c\in\R^q\) (\lec{9}, Theorem~3(b)): knowing \(z_1\) can only
reduce the variance.

\Needspace{10\baselineskip}
\begin{example}\label{h-exm-conditional}
For \(z\) of Example~\ref{h-exm-marginal}, find the conditional distribution of \((z_3,z_4)'\) given \(z_1=a_1\),
\(z_2=a_2\).
\end{example}
\textit{Solution.}
\begin{board}[Theorem~\ref{h-thm-conditional} with \((z_1,z_2)'\), \((z_3,z_4)'\) as its \(z_1\), \(z_2\), \(p=q=2\); the blocks of \(\Sigma\) (Example~\ref{h-exm-marginal}); \(x'\Sigma_{11}x=4x_1^2+x_2^2>0\)
for \(x\ne0\): \(\Sigma_{11}\) is PD:]
\boardtop
\begin{align*}
\Sigma_{11}&=\diag(4,1),\qquad\Sigma_{11}^{-1}=\diag(\tfrac14,1),\qquad
\Sigma_{21}=\begin{pmatrix}4&2\\0&2\end{pmatrix},\qquad\Sigma_{21}\Sigma_{11}^{-1}=\begin{pmatrix}1&2\\0&2\end{pmatrix}\\
\Sigma_{22\cdot1}&=\begin{pmatrix}9&6\\6&9\end{pmatrix}-\begin{pmatrix}1&2\\0&2\end{pmatrix}\begin{pmatrix}4&0\\2&2\end{pmatrix}
=\begin{pmatrix}9&6\\6&9\end{pmatrix}-\begin{pmatrix}8&4\\4&4\end{pmatrix}=\begin{pmatrix}1&2\\2&5\end{pmatrix}
\end{align*}
\end{board}
\begin{board}[The conditional mean (Theorem~\ref{h-thm-conditional}), \(\mu_1=(1,0)'\), \(\mu_2=(2,2)'\):]
\[
\begin{pmatrix}2\\2\end{pmatrix}+\begin{pmatrix}1&2\\0&2\end{pmatrix}\begin{pmatrix}a_1-1\\a_2\end{pmatrix}
=\begin{pmatrix}a_1+2a_2+1\\2a_2+2\end{pmatrix},\qquad
N_2\left(\begin{pmatrix}a_1+2a_2+1\\2a_2+2\end{pmatrix},\begin{pmatrix}1&2\\2&5\end{pmatrix}\right)
\]
The first entry agrees with Example~\ref{h-exm-indep}: \(z_3=z_1+2z_2+e\), \(e\sim N(1,1)\) independent of
\((z_1,z_2)'\). Knowing \(z_1\), \(z_2\) reduces the variances of \(z_3\), \(z_4\) from \(9\), \(9\) to \(1\), \(5\).
\end{board}

\Needspace{19\baselineskip}
\subsection*{4. The chi-square and \(t\) distributions}

\begin{proposition}\label{h-prop-chisq-mean}
If \(u\sim\chi^2_r\), then \(\expc(u)=r\).
\end{proposition}
\begin{proof}
\(u\) has the distribution of \(w_1^2+\dots+w_r^2\) (\ref{f-chisq}), and \(\expc(w_i^2)=\var(w_i)+\expc(w_i)^2=1\)
(Proposition~\ref{h-prop-moments} with \(n=1\); \lec{9}, Theorem~4(d)); so
\(\expc(u)=\expc(w_1^2)+\dots+\expc(w_r^2)=r\) (\ref{f-prob}).
\end{proof}

\Needspace{16\baselineskip}
\begin{theorem}\label{h-thm-t}
Let \(T\sim t_\nu\), \(\alpha\in(0,1)\), \(c\in\R\). Then \(-T\sim t_\nu\), \(\prb(T<-c)=\prb(T>c)\),
\[
\prb(T<-t_{\nu,1-\alpha})=\alpha,\qquad\prb\bigl(-t_{\nu,1-\alpha/2}\leqslant T\leqslant t_{\nu,1-\alpha/2}\bigr)=1-\alpha,
\qquad t_{\nu,1-\alpha/2}>0
\]
\end{theorem}
\begin{proof}
\begin{board}[\(T=Z/\sqrt{V/\nu}\), \(Z\sim N(0,1)\), \(V\sim\chi^2_\nu\) independent (\ref{f-chisq}); \(-Z\sim N(0,1)\)
(Theorem~\ref{h-thm-linear}, \(A=-1\)), independent of \(V\) (\ref{f-functions}); \(k=t_{\nu,1-\alpha/2}\):]
\boardtop
\begin{align*}
-T&=(-Z)/\sqrt{V/\nu}\sim t_\nu&&\text{(\ref{f-chisq})}\\
\prb(T<-c)&=\prb(-T>c)=\prb(T>c)&&\text{(\ref{f-prob}: same distribution)}\\
\prb(T<-t_{\nu,1-\alpha})&=\prb(T>t_{\nu,1-\alpha})=\alpha&&\text{(\ref{f-chisq})}\\
\prb(-k\leqslant T\leqslant k)&=1-\prb(T<-k)-\prb(T>k)=1-\tfrac\alpha2-\tfrac\alpha2&&\text{(line 2 with \(c=k\); \ref{f-prob}, \ref{f-chisq})}\\
\prb(T<0)&=\prb(T>0),\quad\prb(T=0)=0\implies\prb(T\leqslant0)=\tfrac12&&\text{(line 2 with \(c=0\); \ref{f-prob}, \ref{f-chisq})}\\
\prb(T\leqslant0)&=\tfrac12<1-\tfrac\alpha2=\prb(T\leqslant k)\implies k>0&&\text{(\ref{f-chisq}: increasing)}
\end{align*}
\end{board}
\nopagebreak
\end{proof}

\Needspace{19\baselineskip}
\subsection*{5. Standardization and quadratic forms}

\Needspace{16\baselineskip}
\begin{proposition}\label{h-prop-sqrt}
Let \(\Sigma\) be PD, \(\Sigma=Q\Lambda Q'\) (\lec{8}, Theorem~9), so \(\lambda_i>0\) (\lec{8}, Theorem~14(a));
\(\Lambda^{\pm1/2}=\diag(\lambda_1^{\pm1/2},\dots,\lambda_n^{\pm1/2})\),
\(\Sigma^{1/2}=Q\Lambda^{1/2}Q'\), \(\Sigma^{-1/2}=Q\Lambda^{-1/2}Q'\). Then \(\Sigma^{1/2}\), \(\Sigma^{-1/2}\) are
symmetric, and
\[
\Sigma^{1/2}\Sigma^{1/2}=\Sigma,\qquad\Sigma^{1/2}\Sigma^{-1/2}=I,\qquad\Sigma^{-1/2}\Sigma\Sigma^{-1/2}=I,\qquad
\Sigma^{-1/2}\Sigma^{-1/2}=\Sigma^{-1}
\]
\end{proposition}
\begin{proof}
\begin{board}[\(Q'Q=I=QQ'\) (\ref{f-orthonormal}); products of diagonal matrices (\ref{f-views}):]
\boardtop
\begin{align*}
(Q\Lambda^{\pm1/2}Q')'&=Q\Lambda^{\pm1/2}Q'&&\text{(\ref{f-rules}; \(\Lambda^{\pm1/2}\) diagonal)}\\
\Sigma^{1/2}\Sigma^{1/2}&=Q\Lambda^{1/2}Q'Q\Lambda^{1/2}Q'=Q\Lambda Q'=\Sigma\\
\Sigma^{1/2}\Sigma^{-1/2}&=Q\Lambda^{1/2}\Lambda^{-1/2}Q'=QQ'=I\\
\Sigma^{-1/2}\Sigma\Sigma^{-1/2}&=Q\Lambda^{-1/2}\Lambda\Lambda^{-1/2}Q'=QQ'=I\\
(\Sigma^{-1/2}\Sigma^{-1/2})\Sigma&=Q\Lambda^{-1}Q'Q\Lambda Q'=QQ'=I\implies\Sigma^{-1/2}\Sigma^{-1/2}=\Sigma^{-1}
&&\text{(\ref{f-inverse})}
\end{align*}
\end{board}
\nopagebreak
\end{proof}

\Needspace{12\baselineskip}
\begin{theorem}\label{h-thm-standard}
Let \(z\sim N_n(\mu,\Sigma)\), \(\Sigma\) PD, and \(w=\Sigma^{-1/2}(z-\mu)\). Then \(w\sim N_n(0,I_n)\): the \(w_i\)
are independent and \(N(0,1)\); \(z=\mu+\Sigma^{1/2}w\); and \((z-\mu)'\Sigma^{-1}(z-\mu)=w'w\sim\chi^2_n\).
\end{theorem}
\begin{proof}
\begin{board}[Theorem~\ref{h-thm-linear} with \(A=\Sigma^{-1/2}\), \(b=-\Sigma^{-1/2}\mu\); \(\Sigma^{-1/2}\) is symmetric
(Proposition~\ref{h-prop-sqrt}):]
\boardtop
\begin{align*}
w&\sim N_n\bigl(\Sigma^{-1/2}\mu-\Sigma^{-1/2}\mu,\ \Sigma^{-1/2}\Sigma\Sigma^{-1/2}\bigr)=N_n(0,I_n)
&&\text{(Proposition~\ref{h-prop-sqrt})}\\
&\implies w_1,\dots,w_n\text{ independent, }N(0,1)&&\text{(Theorem~\ref{h-thm-indep}(b))}\\
\Sigma^{1/2}w&=\Sigma^{1/2}\Sigma^{-1/2}(z-\mu)=z-\mu\\
w'w&=(z-\mu)'\Sigma^{-1/2}\Sigma^{-1/2}(z-\mu)=(z-\mu)'\Sigma^{-1}(z-\mu)\sim\chi^2_n&&\text{(\ref{f-rules}; \ref{f-chisq})}
\end{align*}
\end{board}
\nopagebreak
\end{proof}
So \(\prb\bigl((z-\mu)'\Sigma^{-1}(z-\mu)\leqslant\chi^2_{n,1-\alpha}\bigr)=1-\alpha\) (\ref{f-chisq}).

\Needspace{13\baselineskip}
\begin{theorem}\label{h-thm-quadratic-chisq}
Let \(z\sim N_n(\mu,\sigma^2I_n)\) (\(\sigma>0\)), and \(A\in\R^{n\times n}\) symmetric, \(A^2=A\), \(\rank A=r\geqslant1\),
\(A\mu=0\). Then \(z'Az/\sigma^2\sim\chi^2_r\).
\end{theorem}
\begin{proof}
\begin{board}[\(A=\sum_i\lambda_iq_iq_i'\) (\lec{8}, Theorem~9(a),(d)), \(\lambda_i\in\{0,1\}\) (\lec{8}, Theorem~11),
\(\#\{i:\lambda_i=1\}=\rank A=r\) (\lec{8}, Theorem~9(b)); number the \(q_i\) so that
\(\lambda_1=\dots=\lambda_r=1\); \(Q_1=[q_1\ \cdots\ q_r]\):]
\boardtop
\begin{align*}
A&=q_1q_1'+\dots+q_rq_r'=Q_1Q_1',\qquad Q_1'Q_1=I_r&&\text{(\ref{f-views}; \ref{f-orthonormal})}\\
Q_1'\mu&=Q_1'Q_1Q_1'\mu=Q_1'A\mu=0\\
u&=\tfrac1\sigma Q_1'z\sim N_r\bigl(0,\tfrac1{\sigma^2}Q_1'(\sigma^2I_n)Q_1\bigr)=N_r(0,I_r)&&\text{(Theorem~\ref{h-thm-linear})}\\
z'Az/\sigma^2&=(Q_1'z)'(Q_1'z)/\sigma^2=u'u=u_1^2+\dots+u_r^2\sim\chi^2_r
&&\text{(Theorem~\ref{h-thm-indep}(b); \ref{f-chisq})}
\end{align*}
\end{board}
\nopagebreak
\end{proof}

\Needspace{11\baselineskip}
\begin{theorem}\label{h-thm-quadratic-indep}
Let \(z\sim N_n(\mu,\sigma^2I_n)\), and \(A,B\in\R^{n\times n}\) symmetric with \(A^2=A\), \(B^2=B\).
\begin{enumerate}
\item If \(AB=0\), then \(z'Az\) and \(z'Bz\) are independent.
\item If \(K\in\R^{m\times n}\) and \(KA=0\), then \(Kz\) and \(z'Az\) are independent.
\end{enumerate}
\end{theorem}
\begin{proof}
\begin{board}[(a) \(A'A=A^2=A\): \(z'Az=z'A'Az=(Az)'(Az)\), a function of \(Az\); likewise \(z'Bz=(Bz)'(Bz)\);
Theorem~\ref{h-thm-indep}(c), \ref{f-functions}:]
\[
A(\sigma^2I_n)B'=\sigma^2AB=0\implies Az,\ Bz\text{ independent}\implies z'Az,\ z'Bz\text{ independent}
\]
\end{board}
\begin{board}[(b) \(z'Az=(Az)'(Az)\) as in (a); Theorem~\ref{h-thm-indep}(c), \ref{f-functions}:]
\[
K(\sigma^2I_n)A'=\sigma^2KA=0\implies Kz,\ Az\text{ independent}\implies Kz,\ z'Az\text{ independent}
\]
\end{board}
\nopagebreak
\end{proof}

\Needspace{19\baselineskip}
\subsection*{6. One normal sample}

\Needspace{16\baselineskip}
\begin{theorem}\label{h-thm-one-sample}
Let \(n\geqslant2\), \(x_1,\dots,x_n\) independent \(N(\mu,\sigma^2)\) variables (\(\sigma>0\)),
\(\bar x=\frac1n\sum_ix_i\), \(s^2=\frac1{n-1}\sum_i(x_i-\bar x)^2\) and \(s=\sqrt{s^2}\).
\begin{enumerate}
\item \(\bar x\sim N(\mu,\sigma^2/n)\).
\item \((n-1)s^2/\sigma^2\sim\chi^2_{n-1}\).
\item \(\bar x\) and \(s^2\) are independent.
\item \(\sqrt n(\bar x-\mu)/s\sim t_{n-1}\).
\end{enumerate}
\end{theorem}
\begin{proof}
\(x=(x_1,\dots,x_n)'\sim N_n(\mu\one,\sigma^2I_n)\) (Theorem~\ref{h-thm-indep}(b)), and \((n-1)s^2=x'C_nx\)
(\ref{f-centering}).

(a) Corollary~\ref{h-cor-marginal}(b) with \(a=\frac1n\one\): \(a'x=\bar x\), \(a'(\mu\one)=\mu\),
\(a'(\sigma^2I_n)a=\sigma^2/n\).
\begin{board}[(b) \(C_n'=C_n=C_n^2\), \(C_n(\mu\one)=0\) (\ref{f-centering}), \(\rank C_n=n-1\geqslant1\) (\lec{8},
Example~5):]
\[
\frac{(n-1)s^2}{\sigma^2}=\frac{x'C_nx}{\sigma^2}\sim\chi^2_{n-1}\qquad\text{(Theorem~\ref{h-thm-quadratic-chisq}, \(A=C_n\))}
\]
\end{board}
\begin{board}[(c) \(K=\frac1n\one'\): \(Kx=\bar x\), \(KC_n=\frac1n(C_n\one)'=0\) (\ref{f-rules}, \ref{f-centering}):]
\[
\bar x\text{ and }x'C_nx\text{ independent}\implies\bar x\text{ and }s^2=\tfrac1{n-1}x'C_nx\text{ independent}
\qquad\text{(Theorem~\ref{h-thm-quadratic-indep}(b); \ref{f-functions})}
\]
\end{board}
\begin{board}[(d) \(Z=\sqrt n(\bar x-\mu)/\sigma\sim N(0,1)\) ((a); Theorem~\ref{h-thm-linear}),
\(V=(n-1)s^2/\sigma^2\sim\chi^2_{n-1}\) ((b)), independent ((c); \ref{f-functions}); \(\prb(V=0)=0\) (\ref{f-chisq}),
so \(s>0\) with probability \(1\):]
\[
\frac{\sqrt n(\bar x-\mu)}s=\frac{\sqrt n(\bar x-\mu)/\sigma}{\sqrt{(n-1)s^2/\bigl((n-1)\sigma^2\bigr)}}
=\frac Z{\sqrt{V/(n-1)}}\sim t_{n-1}\qquad\text{(\ref{f-chisq})}
\]
\end{board}
\nopagebreak
\end{proof}

\begin{definition}\label{h-def-confidence}
Let observations be modeled as random variables whose distribution depends on parameters, and \(\alpha\in(0,1)\). An
interval \([L,U]\) whose endpoints \(L\leqslant U\) are computed from the observations is a \textbf{confidence
interval} with level \(1-\alpha\) for a parameter \(\theta\in\R\) if \(\prb(L\leqslant\theta\leqslant U)=1-\alpha\) for
all values of the parameters.
\end{definition}

\Needspace{13\baselineskip}
\begin{example}\label{h-exm-one-sample}
In Theorem~\ref{h-thm-one-sample}, let \(\alpha\in(0,1)\), \(T=\sqrt n(\bar x-\mu)/s\), \(V=(n-1)s^2/\sigma^2\). Find
confidence intervals with level \(1-\alpha\) for \(\mu\) and for \(\sigma^2\); compute them for five measurements
\(2,4,5,6,8\), modelled as independent \(N(\mu,\sigma^2)\) variables; when does \(3\) lie in the interval for \(\mu\)?
\end{example}
\textit{Solution.}
\begin{board}[\(c=t_{n-1,1-\alpha/2}>0\), \(T\sim t_{n-1}\) (Theorems~\ref{h-thm-t}, \ref{h-thm-one-sample}(d));
\(s>0\):]
\boardtop
\begin{align*}
\prb(-c\leqslant T\leqslant c)&=1-\alpha&&\text{(Theorem~\ref{h-thm-t}; all \(\mu\), \(\sigma\))}\\
-c\leqslant\frac{\sqrt n(\bar x-\mu)}s\leqslant c&\iff\bar x-\frac{cs}{\sqrt n}\leqslant\mu\leqslant\bar x+\frac{cs}{\sqrt n}
\end{align*}
So \([\bar x-cs/\sqrt n,\ \bar x+cs/\sqrt n]\) is a confidence interval for \(\mu\) with level \(1-\alpha\)
(Definition~\ref{h-def-confidence}).
\end{board}
\begin{board}[\(c_1=\chi^2_{n-1,\alpha/2}\), \(c_2=\chi^2_{n-1,1-\alpha/2}\), \(V\sim\chi^2_{n-1}\)
(Theorem~\ref{h-thm-one-sample}(b)); \ref{f-chisq}:]
\boardtop
\begin{align*}
\prb(c_1\leqslant V\leqslant c_2)&=\prb(V\leqslant c_2)-\prb(V<c_1)=1-\tfrac\alpha2-\tfrac\alpha2&&(\prb(V=c_1)=0)\\
\prb(V\leqslant0)&=0<\tfrac\alpha2<1-\tfrac\alpha2\implies0<c_1<c_2\\
c_1\leqslant\frac{(n-1)s^2}{\sigma^2}\leqslant c_2&\iff\frac{(n-1)s^2}{c_2}\leqslant\sigma^2\leqslant\frac{(n-1)s^2}{c_1}
\end{align*}
So \([(n-1)s^2/c_2,\ (n-1)s^2/c_1]\) is a confidence interval for \(\sigma^2\) with level \(1-\alpha\).
\end{board}
\begin{board}[The data, \(n=5\):]
\boardtop
\begin{align*}
\bar x&=\tfrac{25}5=5,\qquad x_i-\bar x:\ -3,-1,0,1,3,\qquad s^2=\tfrac{9+1+0+1+9}4=5\\
\tfrac s{\sqrt5}&=1,\qquad(n-1)s^2=20\\
\mu&:\ \bigl[5-t_{4,1-\alpha/2},\ 5+t_{4,1-\alpha/2}\bigr],\qquad\sigma^2:\ \bigl[20/\chi^2_{4,1-\alpha/2},\ 20/\chi^2_{4,\alpha/2}\bigr]\\
\mu=3&:\ \frac{\sqrt5(\bar x-3)}s=\frac{\sqrt5\cdot2}{\sqrt5}=2,\qquad3\in\bigl[5-t_{4,1-\alpha/2},\ 5+t_{4,1-\alpha/2}\bigr]
\iff t_{4,1-\alpha/2}\geqslant2
\end{align*}
\end{board}

\bigskip
\textbf{Practice.} The Exercises and Problems of the textbook, Chapters~10 and~11: normal vectors, linear
transformations and marginal distributions, independence, conditional distributions, the chi-square, \(t\) and \(F\)
distributions, quadratic forms in normal vectors, and one normal sample.

\end{document}
